\par\noindent\SourceLeaf{59}%
or more briefly
\[
= \frac{1}{2^{2h}}\,{}^{\nu}\mathfrak{B}^{(2h)} .  {}^{2h}\mathfrak{B}^{(h)}\,\frac{\nu+2}{h+1}\,n^{2h+1}.
\]
Next, taking $n$ as variable, form $\int An\,\mathrm{d}n$, whose initial term is $=\frac{1}{2}nn$ and whose general term $h^{\text{tus}}$ is
\[
= \frac{1}{2^{2h}}\,{}^{\nu}\mathfrak{B}^{(2h)} .  {}^{2h}\mathfrak{B}^{(h)}\,\frac{1}{2h+2}\,n^{2h+2}.
\]
This new series, if multiplied by
\[
\frac{2(\nu+2)}{n}
\]
will give the initial term $=(\nu+2)n$ and the general term $h^{\text{tum}}$
\[
= \frac{1}{2^{2h}}\,{}^{\nu}\mathfrak{B}^{(2h)} .  {}^{2h}\mathfrak{B}^{(h)}\,\frac{\nu+2}{h+1}\,n^{2h+1}.
\]
Since the terms on both sides agree, we obtain
\[
2n .  A + A^{(1)} = \frac{2 . (\nu+2)}{n}\int An\,\mathrm{d}n,
\]
with the integral taken without an additive constant. It follows that
\[
A^{(1)} = \frac{2 . (\nu+2)}{n}\int An\,\mathrm{d}n - 2n .  A.
\]
\noindent\hspace*{\parindent}Thus this formula reduces $A^{(1)}$ to the already found $A$, and only the latter remains to be derived from the formation of the series itself, as was done before.

8. Having accomplished what could be hoped for in posing the problem in its most general form, we now turn to special cases, in which the general formulas can sometimes be simplified. An example is when $\nu$ is a positive integer. For then each of the series previously found for expressing the coefficients of $(1+n\cos\varphi)^{\nu}$ must terminate, since each contains, in its own order, the binomial coefficients belonging to the exponent $\nu$. Hence the expression for each coefficient becomes finite. Similarly, when $\nu$ is a negative integer, we can obtain finite expressions for the individual coefficients. Since the formula is chiefly used in this case, it is worth spending some time working it out. The procedure needed for this investigation is carried out in two parts. The first is to find finite expressions for the coefficients when $\nu=-1$. The second is to show how, given the coefficients of $(1+n\cos\varphi)^{\nu}$, those of the power one degree lower,
